\section*{Reading these results}
All forecasts concern monthly log realized variance, log(RV + 1e-8). The common test period is July 2006 to December 2024 and contains 222 months. The pooled sample contains 153 factors and the market portfolio. The raw sample begins in 1986. Expanding preprocessing requires 36 initial observations. The initial estimation split contains 119 training months and 90 validation months, followed by a 209-month refit. Training then expands, validation remains fixed at 90 months, and models are selected each June.
The market portfolio is estimated both within the pooled factor panel and as a separate time series.
Pooled RMSE is the square root of the mean squared error across all series and months. It is therefore distinct from the average factor RMSE. R-squared is computed relative to the mean within each evaluation sample.
EW gives HARX, ENET, RF, XGB, NN1, NN2 and NN3 one-seventh each within each target-month, on the log scale. Each NN architecture counts once despite its ten internal members. Exponentiating this average is not the arithmetic mean of variance-level forecasts. No missing-member renormalization is allowed.
Validation scores describe selected hyperparameters on overlapping 90-month windows. Neural network tuning evaluates one network, while test forecasts average ten refitted networks. The comparison may also reflect refitting, window length, and changes in volatility conditions.
Feature summaries are descriptive linear surrogate loadings. Gross and net shares exclude the intercept. Exact collinearity can produce zero loadings and does not establish economic irrelevance. Predictor scales differ, and the equal-weighted forecast has no fitted surrogate model.
The sequential RMSE differences are arithmetic comparisons across model classes. They should not be interpreted as causal contributions. Elastic net adds shrinkage, while nonlinear learners also change regularization and estimation.
